% Eigenständiger Prosabeweis; Hauptsatz in B28, sechs Hilfsresultate in den Beweistabellen.
\section*{Warum darf man die Klammern weglassen?}

Bei einem Produkt aus vier Faktoren schreiben wir häufig einfach
\(a\star b\star c\star d\). Eine binäre Operation verarbeitet jedoch
immer nur zwei Eingaben. Um tatsächlich zu rechnen, müssen wir festlegen,
welche beiden Teilergebnisse jeweils miteinander verknüpft werden.
Bei unveränderter Reihenfolge der vier Faktoren gibt es fünf vollständige
Klammerungen:
\[
  \begin{aligned}
    &((a\star b)\star c)\star d,\\
    &(a\star(b\star c))\star d,\\
    &(a\star b)\star(c\star d),\\
    &a\star((b\star c)\star d),\\
    &a\star(b\star(c\star d)).
  \end{aligned}
\]
Warum ist diese Liste vollständig? Die äußerste Operation trennt die
Faktorenfolge nach dem ersten, zweiten oder dritten Faktor. Im ersten
Fall bleiben rechts drei Faktoren mit zwei möglichen Klammerungen;
im zweiten Fall haben beide Zweierblöcke jeweils nur eine Klammerung;
im dritten Fall bleiben links wieder zwei Möglichkeiten. So entstehen
genau die fünf angegebenen Formen. Die Faktoren dürfen dabei gleiche
Werte haben; verschieden sind zunächst die Klammerungen.

Dass die Werte übereinstimmen, folgt nicht schon daraus, dass eine
Operation gegeben ist. Für die Subtraktion auf \(\mathbb Z\) ergibt
dieselbe Faktorenfolge \(8,3,2,1\) zum Beispiel
\[
  ((8-3)-2)-1=2,
  \qquad
  (8-3)-(2-1)=4.
\]
Beide Rechnungen sind wohldefiniert, aber ihre Ergebnisse sind verschieden.
Eine solche Operation erlaubt also nicht, die Klammern ohne weitere
Festlegung fortzulassen.

In einer Halbgruppe gilt dagegen das Assoziativgesetz
\[
  (x\star y)\star z=x\star(y\star z)
  \qquad(x,y,z\in A).
\]
Es spricht zunächst über drei Elemente. Diese Elemente dürfen selbst
Ergebnisse früherer Rechnungen sein. Mit \(x=a\star b\), \(y=c\)
und \(z=d\) erhalten wir beispielsweise
\[
  ((a\star b)\star c)\star d=(a\star b)\star(c\star d).
\]
Das erklärt einen Übergang in der Liste. Zu beweisen bleibt, dass
\emph{jede} vollständige Klammerung einer beliebig langen, endlichen
Faktorenfolge denselben Wert liefert. Eine Liste von Einzelrechnungen
für vier oder fünf Faktoren genügt dafür noch nicht.

\par\Needspace{12\baselineskip}
\section*{Aussage und Beweisidee}

Wir setzen zunächst eine Menge \(A\) und eine binäre Operation
\(\star\colon A\times A\to A\) voraus. Die Schreibweise
\(\mathsf{BinOp}(A,\star)\) aus Band~27 besagt genau, dass zu
jedem geordneten Paar von Elementen aus \(A\) ein eindeutig bestimmter
Wert in \(A\) gehört. Dadurch können wir auch Zwischenergebnisse wieder
verknüpfen. Assoziativität ist eine zusätzliche Eigenschaft dieser
Operation. Eine Halbgruppe \((A,\star)\) erfüllt beide Voraussetzungen.

\medskip\noindent\textbf{Satz über die Klammerungsunabhängigkeit.}
Sei \((A,\star)\) eine Halbgruppe. Für jede endliche, nichtleere
Folge von Faktoren aus \(A\) haben alle vollständigen binären
Klammerungen, die die Reihenfolge und sämtliche Vorkommen der Faktoren
beibehalten, denselben Wert. Dieser Wert ist das vollständig von links
geklammerte Produkt. Bei nur einem Faktor ist der Wert dieser Faktor selbst.

Dies ist der Inhalt von
\FormulaRefAuto{SemigroupBracketingIndependence}[thm] in Band~28.
Die Einschränkung auf nichtleere Folgen gehört zum Satz: Ein neutrales
Element wird für eine Halbgruppe nicht vorausgesetzt.

Der Beweis vergleicht nicht jede Klammerung einzeln mit jeder anderen.
Stattdessen vergleichen wir jede mit einer fest gewählten Form: Wir
beginnen beim ersten Faktor und verknüpfen das bisherige Ergebnis
jeweils mit dem nächsten Faktor. Diese \emph{Linksfaltung} liefert
für unsere vier Faktoren den Wert
\[
  ((a\star b)\star c)\star d.
\]
Wir zeigen zuerst, dass man eine Faktorenfolge in zwei nichtleere
aufeinanderfolgende Blöcke zerlegen, beide Blöcke für sich links falten
und ihre Werte anschließend verknüpfen darf. Dieses \emph{Blockgesetz}
macht die Assoziativität für längere Folgen nutzbar.

Danach betrachten wir den letzten Rechenschritt einer beliebigen
Klammerung. Er verknüpft den Wert einer linken Teilklammerung mit dem
Wert einer rechten Teilklammerung. Sind beide Werte bereits als
Linksfaltungen ihrer Faktorenfolgen bekannt, liefert das Blockgesetz
auch für die ganze Klammerung die Linksfaltung. Eine Induktion über den
Aufbau der Klammerungen führt diesen Gedanken für alle Fälle aus.

\section*{Faktorenfolge, Klammerungsbaum und Wert}

Für den Beweis müssen wir drei Dinge auseinanderhalten: die geordnete
Folge der Faktoren, den Plan ihrer Verknüpfung und das daraus berechnete
Element von \(A\). Band~27 stellt die dafür benötigten Wörter,
Klammerungsbäume sowie Induktions- und Rekursionssätze bereit.
Wir erklären hier die verwendeten Begriffe und Eigenschaften. Ihre
mengentheoretischen Konstruktionen werden in diesem Beweis vorausgesetzt.

\medskip\noindent\textbf{Wörter halten die Reihenfolge fest.}
Ein nichtleeres endliches Wort über \(A\) ist eine endliche Folge
\(w=\langle a_1,\ldots,a_n\rangle\) mit \(n\geq1\) und
\(a_i\in A\). Wiederholungen sind erlaubt und bleiben als einzelne
Vorkommen erhalten. Das Wort \(\langle a,a,b\rangle\) ist daher
weder eine Menge \(\{a,b\}\) noch schon ein Produkt. Die Menge
aller nichtleeren endlichen Wörter schreiben wir
\(\NonemptyWords{A}\). Das Einbuchstabenwort zu \(a\) ist
\(\LetterWord{a}\).

Die \emph{Konkatenation} \(u\frown v\) hängt die Folge \(v\)
an die Folge \(u\) an. Zum Beispiel gilt
\[
  \langle a,b\rangle\frown\langle c,d\rangle
    =\langle a,b,c,d\rangle.
\]
Beim Aneinanderhängen wird noch kein \(\star\)-Produkt berechnet.
Die Konkatenation ist assoziativ:
\[
  (u\frown v)\frown w=u\frown(v\frown w).
\]
Beide Seiten bestehen aus den Einträgen von \(u\), danach denen
von \(v\), danach denen von \(w\). Diese Aussage über Folgen ist
\FormulaRefAuto{WordConcatenationAssociative}[thm] in Band~27 und
gilt unabhängig von einer Operation auf \(A\). Sie setzt die hier
zu beweisende Klammerungsunabhängigkeit von Produkten nicht voraus.

\medskip\noindent\textbf{Die Linksfaltung wählt einen Rechenplan.}
Die Linksfaltung ist die Funktion
\(\Pi_{\star}\colon\NonemptyWords{A}\to A\), die durch
\[
  \WordFold{\star}{\LetterWord{a}}=a,
  \qquad
  \WordFold{\star}{u\frown\LetterWord{a}}
    =\WordFold{\star}{u}\star a
\]
bestimmt wird. In der zweiten Gleichung sind \(u\) nichtleer und
\(a\in A\). Die Existenz und Eindeutigkeit einer Funktion mit
diesen Eigenschaften sind durch den Wortrekursionssatz gesichert:
In \FormulaRefAuto{NonemptyWordRecursion}[thm] nehmen wir als
Anfangsabbildung \(a\mapsto a\) und als Fortsetzung
\((x,a)\mapsto x\star a\). Beide sind Funktionen mit den
geforderten Wertebereichen, weil \(\star\) eine binäre Operation
auf \(A\) ist. Dies liefert
\FormulaRefAuto{LeftWordFoldExistence}[thm]. Die Rekursionsgleichungen
werden also nicht lediglich als Rechenvorschrift hingeschrieben;
der Satz begründet, dass sie eine einzige Funktion auf allen
nichtleeren Wörtern festlegen.

Für das wiederkehrende Beispiel ergibt die Vorschrift
\[
  \begin{aligned}
    \WordFold{\star}{\langle a,b\rangle}&=a\star b,\\
    \WordFold{\star}{\langle a,b,c\rangle}&=(a\star b)\star c,\\
    \WordFold{\star}{\langle a,b,c,d\rangle}
      &=((a\star b)\star c)\star d.
  \end{aligned}
\]
Hier wurde nur die festgelegte Rechenreihenfolge ausgeführt.
Assoziativität ist für die Definition der Linksfaltung nicht erforderlich;
auch die Subtraktion besitzt eine solche Linksfaltung.

\medskip\noindent\textbf{Bäume halten die Klammerung fest.}
Eine vollständige binäre Klammerung lässt sich als endlicher Baum
darstellen. Jedes Blatt trägt genau einen Faktor. Jeder innere Knoten
besitzt einen linken und einen rechten Teilbaum und steht für deren
Verknüpfung. Links und rechts sind dabei festgelegt. Die Blätter werden
in dieser Reihenfolge gelesen; ein Tausch der beiden Teilbäume würde
im Allgemeinen auch die Faktorenfolge verändern.

Die Menge dieser Bäume heißt \(\BracketTreeSet{A}\).
Ein einzelnes Blatt mit Beschriftung \(a\) schreiben wir
\(\LeafTree{A}{a}\). Aus zwei Bäumen \(S,T\) entsteht
\(\NodeTree{A}{S}{T}\), indem wir sie unter eine neue gemeinsame
Wurzel setzen, \(S\) links und \(T\) rechts. Jeder Baum ist entweder
ein Blatt oder ein solcher Knoten; im Knotenfall sind seine beiden
Teilbäume eindeutig bestimmt. Diese Zerlegung ist eine der in Band~27
begründeten Eigenschaften der Klammerungsbäume.

\BracketingTreesDiagram

\medskip\noindent\textbf{Blattwort und Auswertung lesen denselben Baum verschieden.}
Das \emph{Blattwort} \(\TreeLeafWord{A}{T}\) liest nur die
Beschriftungen von links nach rechts. Es erfüllt
\[
  \begin{aligned}
    \TreeLeafWord{A}{\LeafTree{A}{a}}&=\LetterWord{a},\\
    \TreeLeafWord{A}{\NodeTree{A}{S}{T}}
      &=\TreeLeafWord{A}{S}\frown\TreeLeafWord{A}{T}.
  \end{aligned}
\]
Es ist immer ein nichtleeres Wort. Die Funktion existiert nach dem
Baumrekursionssatz \FormulaRefAuto{FullPlanarBinaryTreeRecursion}[thm]:
Ein Blatt erhält das Einbuchstabenwort, und am Knoten werden die beiden
bereits gewonnenen Wörter aneinandergehängt. Weil die Konkatenation
zweier nichtleerer Wörter wieder nichtleer ist, bleibt diese Rechnung
im vorgesehenen Wertebereich. Die konkrete Anwendung steht in
\FormulaRefAuto{TreeLeafWordExistenceUnique}[thm].

Die \emph{Auswertung} \(\TreeEvaluation{\star}{T}\) berechnet
hingegen ein Element von \(A\):
\[
  \begin{aligned}
    \TreeEvaluation{\star}{\LeafTree{A}{a}}&=a,\\
    \TreeEvaluation{\star}{\NodeTree{A}{S}{T}}
      &=\TreeEvaluation{\star}{S}\star\TreeEvaluation{\star}{T}.
  \end{aligned}
\]
Auch sie ist eine Anwendung desselben Baumrekursionssatzes, nun mit
den Blattwerten \(a\) und der Knotenoperation \(\star\).
\FormulaRefAuto{TreeEvaluationExistenceUnique}[thm] sichert Existenz
und Eindeutigkeit. Wieder genügt dafür
\(\mathsf{BinOp}(A,\star)\). Ein fest gegebener Baum lässt sich
eindeutig auswerten, selbst wenn verschiedene Bäume verschiedene Werte
liefern, wie im Beispiel der Subtraktion.

Ein Baum ist eine Klammerung des Wortes \(w\), wenn sein Blattwort
genau \(w\) ist. Die Menge dieser Bäume ist
\[
  \WordBracketings{A}{w}
    :=\{T\in\BracketTreeSet{A}\mid\TreeLeafWord{A}{T}=w\}.
\]
Dies ist \FormulaRefAuto{WordBracketingSetDef}[def]. Für jedes
nichtleere Wort existiert mindestens eine solche Klammerung, etwa
die vollständig linke. Ihre Existenz ist in
\FormulaRefAuto{WordBracketingExistence}[thm] bewiesen. Der kommende
Satz vergleicht also tatsächlich vorhandene Klammerungen.

\section*{Das Blockgesetz: zwei Teile getrennt auswerten}
\hypertarget{bracketing.reading.block}{}

Von jetzt an sei \((A,\star)\) eine Halbgruppe. Für zwei beliebige
nichtleere Wörter \(u,v\) wollen wir zeigen:
\begin{equation}
  \WordFold{\star}{u\frown v}
    =\WordFold{\star}{u}\star\WordFold{\star}{v}.
  \tag{B}\label{bracketing-reading-block}
\end{equation}
Links werden die beiden Wörter zuerst aneinandergehängt und dann
gemeinsam links gefaltet. Rechts werden beide Blöcke zunächst getrennt
links gefaltet; erst danach werden ihre beiden Werte verknüpft.
Das ist \FormulaRefAuto{SemigroupWordBlockLaw}[thm] in den Beweistabellen.

Für \(u=\langle a,b\rangle\) und \(v=\langle c,d\rangle\)
lautet die Behauptung
\[
  ((a\star b)\star c)\star d=(a\star b)\star(c\star d).
\]
Die Trennstelle liegt hier zwischen \(b\) und \(c\). Das
Blockgesetz erlaubt jede Trennstelle, an der auf beiden Seiten
mindestens ein Faktor verbleibt.

\begin{proof}[Beweis]
Wir halten \(u\in\NonemptyWords{A}\) fest und führen eine
Induktion über das rechte Wort \(v\). Das in Band~27 bewiesene
Prinzip \FormulaRefAuto{NonemptyWordInduction}[thm] verlangt zwei
Schritte: die Aussage für jedes Einbuchstabenwort und den Übergang
von einem nichtleeren Wort zum Wort mit einem rechts angefügten
Buchstaben. Es erfasst alle nichtleeren endlichen Wörter.

\medskip\noindent\textbf{Ein Buchstabe als rechter Block.}
Sei \(v=\LetterWord{x}\) mit \(x\in A\). Aus den beiden
Rekursionsgleichungen der Linksfaltung folgt
\[
  \WordFold{\star}{u\frown\LetterWord{x}}
    =\WordFold{\star}{u}\star x
    =\WordFold{\star}{u}\star\WordFold{\star}{\LetterWord{x}}.
\]
Damit gilt das Blockgesetz für jeden rechten Block der Länge eins.
Dieser Anfang benötigt die Assoziativität noch nicht.

\medskip\noindent\textbf{Ein weiterer Buchstabe.}
Sei \(v\) ein nichtleeres Wort, für das bereits
\[
  \WordFold{\star}{u\frown v}
    =\WordFold{\star}{u}\star\WordFold{\star}{v}
\]
gilt. Wir hängen ein beliebiges \(x\in A\) rechts an \(v\) an.
Um die Rekursionsregel benutzen zu können, schreiben wir das neue
Gesamtwort zunächst in der Form
\[
  u\frown(v\frown\LetterWord{x})
    =(u\frown v)\frown\LetterWord{x}.
\]
Dies ist die schon bekannte Assoziativität des Aneinanderhängens von
Wörtern. Danach rechnen wir
\[
  \begin{aligned}
    \WordFold{\star}{u\frown(v\frown\LetterWord{x})}
      &=\WordFold{\star}{(u\frown v)\frown\LetterWord{x}}\\
      &=\WordFold{\star}{u\frown v}\star x\\
      &=\bigl(\WordFold{\star}{u}\star\WordFold{\star}{v}\bigr)
          \star x\\
      &=\WordFold{\star}{u}\star
          \bigl(\WordFold{\star}{v}\star x\bigr)\\
      &=\WordFold{\star}{u}\star
          \WordFold{\star}{v\frown\LetterWord{x}}.
  \end{aligned}
\]
Die zweite Gleichheit ist die Rekursionsregel für die Linksfaltung.
Die dritte setzt die Induktionsannahme ein. In der vierten verwenden
wir das Assoziativgesetz für die drei Elemente
\(\WordFold{\star}{u}\), \(\WordFold{\star}{v}\) und \(x\).
Sie liegen sämtlich in \(A\), weil die Linksfaltung ihre Werte
in \(A\) annimmt. Die letzte Gleichheit benutzt die Rekursionsregel
nochmals, nun für den rechten Block.

Damit ist der Induktionsschritt bewiesen. Die Wortinduktion liefert
\eqref{bracketing-reading-block} für jedes nichtleere \(v\).
Da auch das anfangs festgehaltene \(u\) beliebig war, gilt das
Blockgesetz für alle \(u,v\in\NonemptyWords{A}\).
\end{proof}

Genau in der beschriebenen Umklammerung der drei Werte geht die
Assoziativität von \(\star\) in diesen Beweis ein. Bei der
Subtraktion würde aus \((p-q)-x\) der im Allgemeinen andere Wert
\(p-(q-x)\). Das erklärt, weshalb das Blockgesetz dort scheitert.
Die zugehörige formale Ableitung steht unter
\BracketingProofLink{SemigroupWordBlockLaw}.

\section*{Die Normalform: jeder Baum hat den Wert seines Blattwortes}
\hypertarget{bracketing.reading.normal}{}

Nun zeigen wir für jeden Klammerungsbaum \(T\):
\begin{equation}
  \TreeEvaluation{\star}{T}
    =\WordFold{\star}{\TreeLeafWord{A}{T}}.
  \tag{N}\label{bracketing-reading-normal}
\end{equation}
Die linke Seite folgt dem im Baum gespeicherten Rechenplan.
Die rechte Seite liest zuerst die Blätter als Wort und verwendet dann
die feste linke Klammerung. Die Gleichheit besagt, dass dieser Wechsel
des Rechenplans den Wert nicht verändert. Sie ist die
\emph{Baum-Normalform} aus
\FormulaRefAuto{SemigroupTreeNormalForm}[thm] in den Beweistabellen.

\begin{proof}[Beweis]
Wir benutzen die Strukturinduktion für Klammerungsbäume aus
\FormulaRefAuto{FullPlanarBinaryTreeStructuralInduction}[thm].
Dazu zeigen wir die Aussage zunächst für jedes einzelne Blatt.
Danach zeigen wir: Gilt sie für zwei Bäume, so gilt sie auch für den
Baum mit diesen beiden Teilbäumen unter einer neuen Wurzel.
Dieser Induktionssatz erfasst alle endlichen Klammerungsbäume.

\par\Needspace{8\baselineskip}
\medskip\noindent\textbf{Ein einzelnes Blatt.}
Sei \(T=\LeafTree{A}{a}\) mit \(a\in A\). Die Auswertung
dieses Baumes ist \(a\). Sein Blattwort ist \(\LetterWord{a}\),
dessen Linksfaltung ebenfalls \(a\) ist. Somit gilt
\[
  \TreeEvaluation{\star}{\LeafTree{A}{a}}
    =a
    =\WordFold{\star}{\LetterWord{a}}
    =\WordFold{\star}{\TreeLeafWord{A}{\LeafTree{A}{a}}}.
\]

\medskip\noindent\textbf{Ein Knoten mit zwei Teilbäumen.}
Seien \(S,T\in\BracketTreeSet{A}\), und die behauptete
Gleichheit gelte bereits für beide. Wir setzen
\[
  R:=\NodeTree{A}{S}{T},\qquad
  u:=\TreeLeafWord{A}{S},\qquad
  v:=\TreeLeafWord{A}{T}.
\]
Die Wörter \(u\) und \(v\) sind nichtleer, denn jeder der
beiden Teilbäume besitzt mindestens ein Blatt. Deshalb dürfen wir
das Blockgesetz auf sie anwenden. Die Induktionsannahmen lauten
in dieser Schreibweise
\[
  \TreeEvaluation{\star}{S}=\WordFold{\star}{u},
  \qquad
  \TreeEvaluation{\star}{T}=\WordFold{\star}{v}.
\]
Außerdem ist \(\TreeLeafWord{A}{R}=u\frown v\): Die Blätter
des linken Teilbaums stehen vor denen des rechten. Damit folgt
\[
  \begin{aligned}
    \TreeEvaluation{\star}{R}
      &=\TreeEvaluation{\star}{S}\star\TreeEvaluation{\star}{T}\\
      &=\WordFold{\star}{u}\star\WordFold{\star}{v}\\
      &\overset{\eqref{bracketing-reading-block}}{=}
        \WordFold{\star}{u\frown v}\\
      &=\WordFold{\star}{\TreeLeafWord{A}{R}}.
  \end{aligned}
\]
Die erste Gleichheit ist die Auswertungsregel an einem Knoten.
Die zweite benutzt beide Induktionsannahmen. Die dritte ist genau
das Blockgesetz in umgekehrter Leserichtung; die letzte setzt das
Blattwort der neuen Wurzel ein. Damit gilt die Aussage auch für \(R\).

Blattfall und Knotenschritt sind bewiesen. Die Strukturinduktion
liefert \eqref{bracketing-reading-normal} für jeden
\(T\in\BracketTreeSet{A}\).
\end{proof}

\BracketingNormalFormDiagram

Bei unserem Baum für \((a\star b)\star(c\star d)\) besitzt die
linke Teilklammerung das Blattwort \(\langle a,b\rangle\), die
rechte das Blattwort \(\langle c,d\rangle\). Die beiden
Teilwerte sind \(a\star b\) und \(c\star d\). Der Knotenschritt
führt ihr Produkt auf die Linksfaltung von
\(\langle a,b,c,d\rangle\) zurück, also auf
\(((a\star b)\star c)\star d\). Genau dasselbe Argument gilt,
wenn die Teilbäume selbst sehr viel größer sind. Ihre Größe spielt
im Knotenschritt keine Rolle; benötigt werden nur die beiden
Induktionsannahmen und das Blockgesetz.

Die Normalform behauptet eine Gleichheit der \emph{Werte}.
Der ursprüngliche Baum und der vollständig linke Baum dürfen als
Klammerungsbäume verschieden bleiben. Der Beweis identifiziert sie
nicht miteinander. Die formale Ableitung der Wertgleichheit steht unter
\BracketingProofLink{SemigroupTreeNormalForm}.

\section*{Die Klammerungsunabhängigkeit}
\hypertarget{bracketing.reading.independence}{}

\begin{proof}[Beweis des angekündigten Satzes]
Sei \(w\in\NonemptyWords{A}\) ein beliebiges nichtleeres Wort,
und seien \(S,T\in\WordBracketings{A}{w}\) zwei seiner
Klammerungen. Nach der Definition der Klammerungsmenge gilt
\[
  \TreeLeafWord{A}{S}=w=\TreeLeafWord{A}{T}.
\]
Auf beide Bäume wenden wir die Normalform an:
\[
  \begin{aligned}
    \TreeEvaluation{\star}{S}
      &\overset{\eqref{bracketing-reading-normal}}{=}
        \WordFold{\star}{\TreeLeafWord{A}{S}}\\
      &=\WordFold{\star}{w}\\
      &=\WordFold{\star}{\TreeLeafWord{A}{T}}\\
      &\overset{\eqref{bracketing-reading-normal}}{=}
        \TreeEvaluation{\star}{T}.
  \end{aligned}
\]
Beide Werte sind somit gleich. Weil das Wort und seine beiden
Klammerungen beliebig waren, gilt dies für jede endliche nichtleere
Faktorenfolge und alle ihre vollständigen binären Klammerungen.
\end{proof}

Insbesondere besitzen die fünf Klammerungen vom Anfang denselben
Wert. Man braucht nicht zehn Vergleiche zwischen je zwei der fünf
Formen: Jede hat den Wert
\(\WordFold{\star}{\langle a,b,c,d\rangle}\).
Dasselbe gemeinsame Vergleichsobjekt funktioniert für jede weitere
endliche Länge. Hier liegt der Nutzen der Normalform.

In Band~28 ist dieser Schluss als
\FormulaRefAuto{SemigroupBracketingIndependence}[thm] festgehalten; seine
tabellarische Ableitung steht unter
\BracketingProofLink{SemigroupBracketingIndependence}.

\clearpage
\section*{Was das Ergebnis erlaubt}

\medskip\noindent\textbf{Produkte ohne Klammern und Rechnungen mit Blöcken.}
Für eine feste endliche nichtleere Faktorenfolge dürfen wir nun
\(a_1\star\cdots\star a_n\) schreiben: Gemeint ist ihr gemeinsamer
Wert unter allen vollständigen Klammerungen. Ebenso dürfen wir
zusammenhängende Blöcke vorab auswerten. Der Beweis des Blockgesetzes
zeigt genau, weshalb das Ergebnis dadurch nicht verändert wird.
Diese Freiheit betrifft die Wahl der Klammern und der
aufeinanderfolgenden Teilrechnungen.

\medskip\noindent\textbf{Funktionskomposition als Beispiel.}
Betrachten wir alle Selbstabbildungen einer Menge \(X\).
Die Komposition zweier solcher Funktionen ist wieder eine
Selbstabbildung. Für drei Selbstabbildungen \(f,g,h\) und jedes
\(x\in X\) gilt
\[
  ((f\circ g)\circ h)(x)
    =f(g(h(x)))
    =(f\circ(g\circ h))(x).
\]
Da die beiden Funktionen an jeder Stelle denselben Wert haben,
sind sie gleich. Die Komposition bildet also eine assoziative
binäre Operation. Unser Satz gilt deshalb auch für endliche
nichtleere Folgen von Selbstabbildungen. Für vier Funktionen ist
beispielsweise
\[
  ((f\circ g)\circ h)\circ j
    =(f\circ g)\circ(h\circ j),
\]
und beide Seiten schicken \(x\) auf \(f(g(h(j(x))))\).
Das ist dieselbe Vierfaktorenrechnung wie zuvor, nun mit Funktionen
als Faktoren. Nach der üblichen Definition der Komposition wird
die rechts stehende Funktion zuerst auf das Argument angewandt.

\medskip\noindent\textbf{Die Faktorenfolge bleibt fest.}
Assoziativität erlaubt keine beliebige Vertauschung. Bei den
Selbstabbildungen von \(\mathbb Z\) mit
\(f(x)=x+1\) und \(g(x)=2x\) ergibt sich etwa
\[
  (f\circ g)(0)=1,
  \qquad
  (g\circ f)(0)=2.
\]
Die Funktionen \(f\circ g\) und \(g\circ f\) sind daher
verschieden. Auch ein mehrfach vorkommender Faktor darf nicht
einfach ausgelassen werden. Unsere Bäume bewahren alle Blätter
und ihre Reihenfolge; genau darauf beruht die Gleichheit ihrer
Blattwörter im letzten Beweisschritt.

\par\Needspace{24\baselineskip}
\medskip\noindent\textbf{Ein Faktor, kein Faktor und unendlich viele Faktoren.}
Bei einem einzigen Faktor besteht der Baum aus einem Blatt, und der
Produktwert ist der Faktor selbst. Für ein leeres Produkt liefert der
Satz hingegen keinen Wert. Besitzt die Halbgruppe zusätzlich ein
neutrales Element \(e\), ist sie ein Monoid; dann kann man das
leere Produkt durch \(e\) ergänzen. Diese zusätzliche Vereinbarung
wird hier nicht benötigt. Ebenso trifft der Satz keine Aussage über
unendliche Produkte. Solche Produkte verlangen weitere Begriffe und
Voraussetzungen.

\par\Needspace{13\baselineskip}
\medskip\noindent\textbf{Die Aufgabe der beiden Induktionen.}
Die Wortinduktion verlängert einen Block um einen letzten Faktor.
An genau diesem Schritt wird die Assoziativität auf drei bereits
vorhandene Werte angewandt. Die Bauminduktion setzt anschließend
zwei Teilklammerungen zusammen und benutzt das bewiesene Blockgesetz.
Sie erweitert damit die feste linke Rechnung auf jede mögliche
vollständige Klammerung. Die Rekursionssätze aus Band~27 haben zuvor
die benötigten Funktionen bereitgestellt; die Induktionen beweisen
nun ihre Beziehungen zueinander.

Der Hauptsatz steht in Band~28. Das Blockgesetz, die Baum-Normalform,
ihre vier Hilfssätze und alle sieben formalen Ableitungen stehen in den
zugehörigen \BracketingProofsLink{}.
Die hier verwendeten Grundlagen über Wörter, Bäume und ihre
Rekursion bleiben in Band~27. Zum Hauptband führt die
\BracketingNotationLink{}.
